\chapter{easy-gl and meta-gl: Profile, Loader, and Dispatch}
\index{easy-gl}\index{meta-gl}\index{OpenGL loader}
\label{ch:easygl-metagl}
\label{ch:easygl-deep-dive}

This book introduced \texttt{easy-gl} as the C\texttt{++}23 OpenGL/OpenGL ES wrapper behind
CNA's five shared-family profiles: \texttt{OPENGLES2}, \texttt{OPENGLES3},
\texttt{OPENGL33}, \texttt{WEBGL1}, and \texttt{WEBGL2}. The public renderer family is really
a two-library stack. \texttt{meta-gl} owns loading, typed calls, context facts, and validation;
\texttt{easy-gl} owns resources, object-oriented state, and RAII.

\section{The ownership split}

\texttt{easy-gl}'s own README frames it as toolkit-independent, meaning independent of a
windowing toolkit like SDL or Qt --- but it is not independent of a lower C\texttt{++} GL
wrapper underneath it. \texttt{easy-gl} itself is built on a sibling project,
\texttt{meta-gl}, which provides the actual procedural, typed, ownership-free OpenGL/GLES
call surface (typed enum-class wrappers, \texttt{std::span}-based views, lightweight typed
handles, no RAII of its own); \texttt{easy-gl} then layers real object-oriented, RAII,
move-only ownership semantics on top. The dependency is real and direct in the build system
(\texttt{meta-gl} is pulled in as a sibling \texttt{add\_subdirectory} and linked publicly),
and in the source: 19 production files include \texttt{meta-gl} headers directly, comprising
three public headers and sixteen implementation files. The layering is worth remembering
precisely because it clarifies what
``toolkit-independent'' actually means here: independent of the host windowing/event toolkit,
not independent of every other C\texttt{++} abstraction below it.

Neither library creates a window or GL context, pumps events, nor presents a frame. The host
owns those operations and supplies a \texttt{GetProcAddress} callback after making its context
current. CNA fills that host role through SDL, while the same libraries remain usable with a
different toolkit.

The size inversion is easy to miss. Easy-gl exposes 30 headers and about 3{,}606 production
code lines; meta-gl exposes only 12 headers but roughly 8{,}856 true C/C\texttt{++} code lines.
The lower procedural layer is about 2.5 times larger because its generated-looking surface is
still a hand-audited typed mapping of hundreds of GL entry points.

\section{Device initialization delegates context truth downward}

\cnaclass{easygl::Device} is the largest single entry point, covering initialization, the
bulk of GL state (clear, viewport/scissor, blend --- including indexed, per-attachment blend
functions for multiple render targets --- depth/stencil, culling, polygon/line mode, sample
coverage), pixel readback (a bounds-checked \texttt{read\_pixels\_robust} alongside the plain
form), the full draw-call family (instanced, indirect, ranged, base-vertex), compute
dispatch, tessellation patch configuration, and debug facilities (labeled debug groups, a
settable debug callback).

Older easy-gl code duplicated version-string parsing. The audited implementation calls
\texttt{metagl::Initialize()} and then copies \texttt{GetContextInfo()} and
\texttt{GetCapabilities()} into its higher-level model. This matters on the web: meta-gl
classifies Emscripten contexts as \texttt{ApiKind::WebGL}, so WebGL no longer masquerades as
ordinary OpenGL ES merely because its shading language is ES-shaped.

Meta-gl also computes a \texttt{DesktopEsTier}: the OpenGL ES tier whose mandatory functions
a desktop OpenGL context can supply. That mapping is what lets CNA's \texttt{OPENGL33}
identity share one implementation with native ES and WebGL profiles without pretending that
their version numbers have identical meanings.

Once detection completes, easy-gl requires vertex arrays, shaders, programs, buffers, and
basic rendering. A missing baseline throws during initialization. Optional entry points ---
queries, samplers, transform feedback, syncs, program pipelines, and texture-level queries ---
are checked through \texttt{metagl::IsFunctionAvailable()} and fail with
\cnaclass{easygl::UnsupportedFeatureException} instead of calling a null function pointer.
WebGL~1 may obtain vertex arrays through \texttt{GL\_OES\_vertex\_array\_object}; separable
program pipelines have no WebGL equivalent, including WebGL~2.

\section{Capability gating: a real per-feature table, not a version-number heuristic}

\texttt{easygl::Capabilities::detect\_common\_features()}%
\index{easygl::Capabilities::detect_common_features()@easygl::Capabilities::\allowbreak{}detect\_common\_features()}
implements a genuine,
per-feature gating table rather than a single ``desktop GL is capability X, GLES is capability
Y'' heuristic, and the two API families frequently diverge in specific, non-obvious ways.
Vertex array objects need GL~3.0 (or the \texttt{GL\_ARB\_vertex\_array\_object} extension) on
desktop, but GLES~3.0 or the differently-named \texttt{GL\_OES\_vertex\_array\_object}
extension on GLES. Framebuffer objects need GL~3.0 or an extension on desktop, but are simply
always present on GLES~2.0 and above, since FBOs are core to ES2 from the start. Tessellation
shaders need GL~4.0 on desktop but only GLES~3.2 --- a case where the GLES requirement is
numerically \emph{lower}. Direct State Access, by contrast, is not modeled for GLES at all,
since it is a desktop-only OpenGL~4.5 capability with no GLES equivalent whatsoever. This
per-feature table --- not a single ``is this GLES'' branch --- is exactly the mechanism
Chapter~\ref{ch:easygl-backend} described CNA's own EasyGL renderer implementation
relying on via \texttt{device.supports(...)} / \texttt{device.require(...)}.

\begin{lstlisting}
// Worked example: querying the per-feature table directly, rather than
// branching on Device::capabilities().is_opengles() and guessing.
if (device.supports(easygl::Feature::TessellationShader))
{
    ConfigurePatchPipeline();
}
else
{
    UseFallbackSubdivision();  // e.g. desktop GL < 4.0, or GLES < 3.2
}

// With the default Config::throw_on_missing_feature == true,
// Device::require() throws immediately if the feature is missing.
device.require(easygl::Feature::FramebufferObject);
\end{lstlisting}

\section{Meta-gl's typed call surface and invalid-input contract}

Meta-gl exports 358 numbered \texttt{metagl::gl*} wrappers. A Python verifier requires the
markers to form one contiguous 1--358 sequence with unique names and cross-checks the mandatory
OpenGL ES 3.0, 3.1, and 3.2 sets (104, 68, and 44 additions) against the vendored Khronos
header. The public type vocabulary contains 105 enum classes in \texttt{Enums.hpp} and 16
struct declarations in \texttt{Types.hpp}: typed object handles, locations, indices, an image
unit, and a bitfield-traits helper. Older continuation counts of 99 enums and 15 structs are
stale.

These handles are intentionally ownership-free:

\begin{lstlisting}
metagl::TextureId texture{rawName};
metagl::glBindTexture(metagl::TextureTarget::Texture2D, texture);

// TextureId does not delete rawName. easygl::Texture owns that lifetime.
\end{lstlisting}

The two layers also choose different failure semantics. Easy-gl reports unsupported features
with C\texttt{++} exceptions. Meta-gl's invalid-input contract calls
\texttt{std::terminate()} even in Release for conditions such as a \texttt{size\_t} that
cannot fit in \texttt{GLsizei}, incomplete matrix data, or an unsupported bitfield. Those are
treated as caller contract violations, not recoverable driver capabilities.

This difference has a build consequence on Emscripten. Easy-gl enables
\texttt{-fwasm-exceptions} globally before adding meta-gl so every object in the combined stack
agrees on the Wasm exception model. A failure-policy choice at the OOP layer therefore affects
how the lower procedural layer must be compiled even though meta-gl itself does not throw for
its input contract.

Meta-gl's small header count also hides unusually strict binary discipline: hidden visibility,
a public export macro, a GNU/Clang version script, SONAME and exported-symbol policy tests, and
an installed-package consumer. It can be found as \texttt{meta-gl 0.3}; an uncommitted
\texttt{0.4.0-snapshot} file in the local worktree is not a released or pinned version.

\section{Real design fixes, and one genuine breaking change worth knowing}

\texttt{easy-gl}'s own planning documents record several real, reasoned design corrections.

\texttt{Program::uniform\_block\_index()} originally returned GL's own raw
\texttt{GL\_INVALID\_INDEX} sentinel on a failed lookup, forcing every caller to know and
compare against a raw GL constant; it now returns \texttt{std::optional<unsigned int>}
instead, so a caller never needs to include a GL header at all just to check for failure.
Several resource classes (\cnaclass{Framebuffer}, \cnaclass{ProgramPipeline},
\cnaclass{TransformFeedback}) were changed from accepting a raw \texttt{unsigned int} GL
handle to accepting a real \cnaclass{Texture} / \cnaclass{Program} reference directly, removing
a whole class of ``passed the wrong handle type by mistake'' bug at the API boundary.

\begin{lstlisting}
// Worked example: both fixes in the same worked snippet. No raw GL
// constant comparison, and no raw GLuint handle to mix up with an
// unrelated resource's handle by mistake.
if (std::optional<unsigned int> blockIndex = shaderProgram.uniform_block_index("Lighting"))
{
    shaderProgram.set_uniform_block_binding(*blockIndex, /* bindingPoint */ 0);
}
// else: no "Lighting" uniform block in this program -- handled without
// ever needing to know or compare against GL_INVALID_INDEX directly.

Framebuffer gbuffer;
gbuffer.attach_texture_2d(FramebufferTarget::Framebuffer, FramebufferAttachment::Depth,
                           TextureTarget::Texture2D, depthTexture, /* level */ 0);
// depthTexture is a real Texture&, not a bare GLuint -- passing a
// Renderbuffer's handle here by mistake is now a compile error, not a
// runtime GL_INVALID_OPERATION discovered later.
\end{lstlisting}

One specific, genuine breaking behavioral change is worth knowing if you are porting code
against an older version of this library: \texttt{Device::clear()} used to call
\texttt{glDisable(GL\_SCISSOR\_TEST)} before every clear, silently overriding whatever scissor
state the caller had already set. This is no longer true --- \texttt{clear()} no longer
touches scissor state at all, and a caller that was relying on the old implicit disable must
now call the scissor-disable method explicitly. The project's own migration notes frame this
as low-impact, since it only matters when scissor testing happens to already be active at the
moment \texttt{clear()} is called --- but it is exactly the kind of silent, implicit
side-effect this book has flagged elsewhere (Chapter~\ref{ch:spritebatch}'s ignored transform
matrix, for instance) as the class of bug most likely to go unnoticed until specifically
tested for.

\subsection{Four small, real RAII additions worth knowing about}

Two small utility classes, added after this chapter's own original pass, are worth a direct
mention because each closes a real, easy-to-get-wrong ergonomic gap rather than adding a new
GL capability. \cnaclass{ScopedDebugGroup} is a two-call RAII pair --- its constructor calls
\texttt{Device::push\_debug\_group}, its destructor unconditionally calls
\texttt{pop\_debug\_group} --- specifically so a labeled debug region in a RenderDoc or NVIDIA
Nsight capture can never be left unbalanced by an early return or an exception unwinding
through the scope:

\begin{lstlisting}
{
    ScopedDebugGroup group(device, "Shadow Pass");
    RenderShadowCasters();
    // pop_debug_group() always runs here, even if RenderShadowCasters()
    // throws or returns early -- no separate try/finally needed.
}
\end{lstlisting}

\cnaclass{UniformCache} addresses the other classic OpenGL ergonomics problem: repeatedly
calling \texttt{glGetUniformLocation} by name is a real, measurable per-draw cost real engines
avoid by caching the resolved integer location once. \texttt{operator[]} looks up a cached
location by name, falling through to a real \texttt{Program::uniform\_location} call and
memoizing the result only on a cache miss --- and, worth knowing explicitly since it is not
stated in the header's own brief comment, a \emph{failed} lookup (a mistyped or since-optimized-out
uniform name, returning $-1$) is memoized exactly the same way a successful one is. Calling
\texttt{operator[]} again with the same wrong name will not re-attempt the GL query and
magically start succeeding once the shader changes --- only an explicit \texttt{invalidate()}
call (needed after any program relink, since old locations are not guaranteed valid against a
newly linked program) clears the cache and allows a fresh lookup:

\begin{lstlisting}
UniformCache uniforms(shaderProgram);
shaderProgram.set_uniform(uniforms["uTime"], elapsedSeconds);
shaderProgram.set_uniform(uniforms["uLightDir"], lightDir.x, lightDir.y, lightDir.z);
// ... later, only after shaderProgram is actually relinked ...
uniforms.invalidate();  // required -- old cached locations are not safe to reuse otherwise
\end{lstlisting}

Two further small RAII utilities, added in the same wave as the two above, round out this
pattern rather than introducing a new one. \cnaclass{ScopedBind} generalizes the same
constructor-runs-first/destructor-always-runs shape into a reusable, resource-agnostic
wrapper --- constructed from any pair of bind/unbind callables, not tied to one specific GL
object type at all:

\begin{lstlisting}
{
    ScopedBind bound(
        [&]{ vao.bind(); },
        [&]{ vao.unbind(); });
    ConfigureVertexAttributes(vao);
    // unbind() always runs here, even on an early return or a thrown
    // exception mid-configuration -- the same guarantee ScopedDebugGroup
    // gives for a debug region, applied generically to any bind/unbind
    // pair -- VertexArray here, but the constructor takes any two
    // callables, not one fixed resource type.
}
\end{lstlisting}

\cnaclass{ResourceRegistration} closes a different, narrower gap: any GL resource that
participates in this library's context-loss recovery mechanism must register itself with a
\cnaclass{ResourceRegistry} to be told when to rebuild itself after a context is lost and
recreated, and must remember to deregister on destruction to avoid the registry holding a
dangling pointer. Before this addition, both calls were the caller's own responsibility to
remember; \cnaclass{ResourceRegistration}'s constructor calls \texttt{registry.add(\&resource)}
and its destructor unconditionally calls \texttt{registry.remove(\&resource)}, the identical
RAII shape as every other addition in this section, applied specifically to the
register/deregister pair a \cnaclass{RecoverableResource} would otherwise have to manage by
hand at every one of its own constructor and destructor sites.

\section{Context loss recovery: two real tiers, one of them built but not yet adopted}

\cnaclass{ResourceRegistration} above closes the register/deregister half of the context-loss
story, but reading \texttt{GenerationTracked.hpp} and the context-lifecycle test suite together
reveals a fuller picture worth spelling out on its own: \texttt{easy-gl} actually has
\emph{two} separate mechanisms for surviving a lost GL context, at two different levels of
commitment, and only one of them is currently used by anything real.

The first tier covers every integer-named GL resource class and is effectively free. Those
eleven classes inherit a common base class,
\texttt{easygl::\allowbreak{}detail::\allowbreak{}GenerationTracked}. It stamps each
resource with the \texttt{meta-gl}-tracked context generation number at creation time
(\texttt{creation\_generation()}) and exposes \texttt{is\_valid\_for\_current\_generation()} so
any caller can cheaply ask ``was this handle created against the context that is actually
current right now, or a since-lost one'' without any registration step at all. The pointer-
shaped \cnaclass{Sync} object is the exception: it has its own creation and no-generation reset
surface, so it can discard a stale handle but cannot answer the generation-validity query.
\texttt{reset\_handle\_no\_gl()} zeroes either representation without issuing a single
\texttt{gl*} call, since a lost context cannot be trusted to process one correctly anyway.

The second tier is the opt-in, full self-healing mechanism this chapter already covers:
\cnaclass{RecoverableResource}'s \texttt{release\_gl\_handle\_only()} / \texttt{recreate\_gl\_resource()}
pair, driven by \cnaclass{ResourceRegistry} (itself a real \texttt{metagl::ContextListener},
receiving \texttt{OnContextLost()} / \texttt{OnContextRestored()} automatically once
\texttt{register\_with\_meta\_gl()} is called). This tier is where the real commitment lives:
implementing it means a resource class must retain enough CPU-side state (raw pixel data, shader
source text, buffer contents) to rebuild the GPU object from scratch, not merely notice that its
handle went stale.

\textbf{A genuinely worth-knowing gap}: the infrastructure has production implementations ---
\cnaclass{ResourceRegistry} dispatches the two lifecycle calls and
\cnaclass{ResourceRegistration} supplies RAII add/remove --- but \emph{zero concrete resource
classes} adopt it. \cnaclass{Texture}, \cnaclass{Buffer}, \cnaclass{VertexArray},
\cnaclass{Framebuffer}, and every other real GL object class inherit \texttt{GenerationTracked}
(tier one) but none of them implement \texttt{RecoverableResource} (tier two). The library's own
planning document confirms this is a known, tracked state rather than an oversight this chapter
is the first to notice: Task U2 (``Test context-loss / \texttt{ResourceRegistry} recovery
cycle'') is explicitly scoped against ``a concrete \texttt{RecoverableResource} subclass''
precisely because no real one currently exists to test against --- and, consistent with that,
\texttt{ContextLifecycleTests.cpp} tests this directly, with two dedicated
\texttt{ResourceRegistry} tests:

\begin{itemize}
  \item \texttt{test\_resource\_registry\_context\_lost}
  \item \texttt{test\_resource\_registry\_wired\_to\_meta\_gl}
\end{itemize}

Both construct a minimal, test-only \texttt{FakeResource} stub rather than exercising any
real resource type, for exactly this reason:

\begin{lstlisting}
// From the real test suite -- note this is a hand-rolled stub, not Texture,
// Buffer, or any other production easy-gl resource class. None of those
// currently implement RecoverableResource.
struct FakeResource : easygl::RecoverableResource
{
    int& released;
    int& recreated;
    void release_gl_handle_only() override { ++released; }
    void recreate_gl_resource()   override { ++recreated; }
};
\end{lstlisting}

The practical upshot for a caller today: a context loss is cheaply \emph{detectable} per
resource via \texttt{is\_valid\_for\_current\_generation()} (tier one, universal), but nothing
in \texttt{easy-gl} itself will \emph{automatically rebuild} a lost \texttt{Texture} or
\texttt{Buffer} for you (tier two, fully built and independently tested as a mechanism, but with
no adopters yet) --- a caller that needs real survive-context-loss behavior today has to detect
staleness via tier one and rebuild the resource manually, or write its own
\cnaclass{RecoverableResource} subclass around its own CPU-side data. This is the same shape of
gap this book has already named on the CNA side more than once (a mechanism exists and is
correctly wired, but nothing yet plugs into it) --- here, one level further down the stack, in
the sibling library CNA's own EasyGL implementation is built on.

\section{Testing without a GPU}

Easy-gl's suite has four hand-rolled, mock-loader binaries: the initialization, resource,
context-lifecycle, and WebGL smoke suites. Their CMake target names all begin
\texttt{easy-gl-} and end in \texttt{-tests}. The mechanism making them useful without any
real GPU or display is worth describing directly: fake \texttt{GLGetProcAddressFn} loaders
return hard-wired callbacks, literal strings such as \texttt{"OpenGL ES 3.0"}, and literal
enum values in place of driver responses. The smallest loaders provide the context-query
handful needed by \texttt{Device::initialize()} and use a non-callable sentinel for other
loaded core functions; the resource suite supplies a broader set of callable stubs for the
operations it actually exercises. This lets context/capability detection and resource
lifecycle assertions execute entirely off-GPU without accidentally issuing a real GL call.

The lifecycle suite drives context loss and restoration; the WebGL suite checks explicit
classification, ES2/ES3 tiers, OES vertex-array gating, and exception-not-crash behavior for
unavailable entry points. The desktop-shaped initialization and resource smoke binaries are
not built or registered under Emscripten because their invented version strings cannot occur
in a real WebGL host. This is an honest test-model boundary, not proof that the corresponding
browser behavior was executed.

Nested build defaults are asymmetric. Easy-gl defaults both tests and examples ON. CNA's
\texttt{web} preset explicitly sets both \texttt{EASYGL\_BUILD\_*} options to OFF, but CNA's
root CMake and renderer-selection helper do not force either value. A direct GL-profile
configuration that does not use that preset or pass its own overrides therefore inherits the
four test binaries, and its SDL hello-triangle example can configure because CNA already supplies
SDL. Meta-gl defaults all tests, GPU tests, examples, documentation, sanitizers, and debug
options OFF when nested. The effective extra work consequently depends on the CNA configure
entry point; the read-only source audit did not build the siblings to measure its wall-clock cost.

Meta-gl's standalone tests add compile, mock-loader, release-contract, thread, desktop-tier,
SONAME, export-symbol, and installed-package checks. A real EGL/Mesa smoke test is separately
opt-in, and requesting GPU tests without ordinary tests is a configure-time error.

\section{Two known, honestly stated limitations}

Two gaps are worth knowing about directly rather than assuming closed.

First, \texttt{Config::enable\_debug\_logging} currently does nothing at all --- no
\texttt{log\_callback} field yet exists on \texttt{Config} to actually receive the output,
despite the flag existing to request it. Second, \cnaclass{Query::result\_u64()} is
implemented using the 32-bit \texttt{glGetQueryObjectuiv} internally, because the underlying
\texttt{meta-gl} layer does not expose the real 64-bit query variant --- meaning it is
suitable for frame counts, but not for a high-resolution GPU timer query, despite its
64-bit-suggesting name. Both are stated plainly in the project's own tracking rather than
silently left for a reader to discover.

\section{The public history was rewritten without changing the trees}

On 2026-08-09, both libraries' public histories were replayed to remove co-author trailers.
The current easy-gl and meta-gl tree hashes are byte-identical to the content CNA had already
accepted, but the old commit objects are unreachable from any branch. Archive tags were
repointed while retaining annotations that name old targets, and the rewritten commits are
unsigned even though the pre-rewrite objects had valid signatures.

For this book, the durable evidence is branch, tree hash, and audit date. The current public
heads are easy-gl \texttt{0b46d35c} and meta-gl \texttt{571d3a62}; older SHAs from the
integration campaign may exist only as dangling local objects and should not be offered as
reader-resolvable citations. The rewrite changed provenance objects, not the implementation
described in this chapter.
